% Generated draft: requires mathematical and source review.
\documentclass[11pt]{article}
\usepackage[margin=25mm]{geometry}
\usepackage{fontspec}
\setmainfont{Latin Modern Roman}
\usepackage{amsmath,amssymb,mathtools,unicode-math}
\setmathfont{Latin Modern Math}
\usepackage{enumitem,needspace,xurl,hyperref}
\setlength{\parindent}{0pt}
\setlength{\parskip}{5pt}
\setlength{\emergencystretch}{3em}
\newcommand{\N}{\mathbb{N}}
\newcommand{\Z}{\mathbb{Z}}
\newcommand{\Q}{\mathbb{Q}}
\newcommand{\R}{\mathbb{R}}
\newcommand{\C}{\mathbb{C}}
\newcommand{\F}{\mathbb{F}}
\newcommand{\A}{\mathbb{A}}
\newcommand{\PP}{\mathbb{P}}
\newcommand{\E}{\mathbb{E}}
\newcommand{\eps}{\varepsilon}
\newcommand{\dd}{\,\mathrm{d}}
\begin{document}
\section*{UnsolvedMath ID: 20001908}
\textbf{Source problem code:} catalogue entry
\subsection{Definitions and mathematical statement}
The extracted Burago--Kleiner Conjecture 7.3.1 asserts that the metric on the fixed \(\Z\)-cover associated with the \(S^1\)-factor of \(M\times S^1\) is a direct product
\[
 (M\times\R,\widetilde g)\simeq(M,g_1)\times(\R,\dd t^2).
\]
The aid interprets the missing preceding setting as a metric without conjugate points on a closed product manifold. The actual preceding hypotheses are not quoted in the supplied record, so the counterexample scope must remain the literal fixed-cover interpretation, rather than an assertion about every possible fuller formulation.

For that interpretation, let \((\Sigma,h)\) be a closed connected oriented hyperbolic surface and let \(\alpha\) be a smooth closed one-form with \([\alpha]\ne0\in H^1_{\mathrm{dR}}(\Sigma)\). On \(\Sigma\times(\R/\Z)\), with \(\dd\theta\) the global angular form, put
\[
 g_\alpha=h+(\dd\theta+\alpha)^2.
\]
Locally \(\alpha=\dd f\), so the coordinate \(u=\theta+f\) makes this a local product. Hence its sectional curvatures are nonpositive and it has no conjugate points. The fixed cyclic cover defined by projection to \(\R/\Z\), equivalently by the kernel of the corresponding homomorphism \(\pi_1(\Sigma\times S^1)\to\Z\), has metric
\[
 \widetilde g_\alpha=h+(\dd t+\alpha)^2\quad\text{on }\Sigma\times\R.
\]
The field \(W=\partial_t\) is unit and parallel, with dual one-form
\(\eta=\dd t+\alpha\). Its periods are the period character
\[
 \pi_1(\Sigma)\longrightarrow\R,\qquad
 [c]\longmapsto\int_c\alpha.
\]
They are nonzero for some loop, so \(\eta\) is not exact.

On the universal cover, a primitive of the lifted \(\alpha\) gives the ordinary product \(\widetilde\Sigma\times\R\). Its hyperbolic factor has no parallel fields, so the only parallel direction is \(W\). If the specified cyclic cover were isometric to \((K,k)\times\R\) with a compact surface factor \(K\), the unit line field of that product would pull back to \(\pm W\), whose dual would be exact. This contradicts the period character. Thus the fixed cover is not such a product, even under an unmarked isometry. In this family the exact criterion is
\[
 (\Sigma\times\R,h+(\dd t+\alpha)^2)
 \text{ is a compact-factor product}\quad\Longleftrightarrow\quad[\alpha]=0.
\]
For the reverse implication, \(\alpha=\dd f\) and \(u=t+f\) give the global product isometry. A product structure on the universal cover is fully consistent with this failure on the specified cyclic cover.

More generally, in a trivialization \(B\times\R\) adapted to a free proper parallel unit translation field \(\partial_t\), the invariant metric has the form
\[
 h_B+(\dd t+A)^2,\qquad \dd A=0.
\]
A product trivialization preserving that line field and the base exists exactly when \([A]=0\). To promote this marked criterion to an obstruction against every unmarked product decomposition, one needs uniqueness of the parallel line, as in the hyperbolic example; if additional parallel directions exist, nonexactness of the chosen dual form alone does not suffice.

The source also suggests broader questions for no-conjugate-points metrics on products and manifolds admitting higher-rank nonpositively curved metrics, without stating them precisely. Its page heading and the start of Section 7.4 about extending conjugate-point notions to magnetic geodesics and natural Hamiltonian systems are adjoining context, not additional conjectures to reconstruct.
\subsection{Short English statement}
A closed nonexact one-form twists a locally product metric on a surface times a circle. The universal cover splits, but the cyclic cover fixed by the original circle projection has a parallel direction with nonzero periods and cannot split with a compact factor. The omitted source hypotheses still need checking.
\subsection{Sources}
\begin{itemize}
\item Permanent UnsolvedMath ID \(20001908\), exact public code catalogue entry; local registry \texttt{aim-geometry-notes.json}, \texttt{notes[245]}, canonical tag conjecture.
\item AIM workshop ``Geodesics,'' Geometry, Section 7.3, Conjecture 7.3.1, attributed in text to Burago--Kleiner: \url{https://aimath.org/pastworkshops/geodesicsproblems.pdf}. No PDF was accessed for this draft. Original literature, proposer/year metadata and full bibliographic citations are absent.
\item Research attempt 1, ``Closed-one-form twists obstruct the asserted cyclic-cover product,'' labelled counterexample, refuted at its run, proved and low novelty confidence; registry solved and SOLVED-IN-LITERATURE. The supplied verification note says preceding-record hypotheses were lost at the boundary but does not quote them in this record. The exact field preserves the extracted statement, not a recovered full hypothesis list.
\end{itemize}
\end{document}
